\documentclass[11pt]{article}

\usepackage[margin=0.8in]{geometry}
\usepackage{amsmath,amssymb}
\usepackage{enumitem}

\setlength{\parindent}{0pt}
\setlength{\parskip}{5pt}

\begin{document}

{\large \textbf{AI-Assisted Program Analysis \& Formal Verification}}

\hfill 

\vspace{8pt}

\textbf{Precondition}

condition which has to be true just before some code
starts executing.

So basically caller has to make sure this is satisfied.

\[
P \quad \rightarrow \quad C
\]

if precondition is not satisfied $\rightarrow$ can get
wrong computation / security problem etc.

One example discussed was Nomad Bridge hack.
The check was incomplete, so an invalid root could pass.

The required condition was

\[
\_root \neq 0
\quad\land\quad
confirmAt[\_root] > 0
\]

so just checking whether the value is non-zero was not
enough.

\vspace{8pt}

\textbf{after the code}

Postcondition = property which should hold immediately
after the code finishes.

Here responsibility is mainly with the callee.

Hoare triple notation:

\[
\{P\}\ C\ \{Q\}
\]

means

\begin{itemize}[leftmargin=18pt,itemsep=1pt]
    \item $P$ : condition before
    \item $C$ : program/code
    \item $Q$ : condition after
\end{itemize}

Example:

\[
\{x>0\}\quad x:=x*2\quad\{x\bmod 2=0\}
\]

The postcondition says the resulting $x$ is even.

\vspace{8pt}

\textbf{strongest postcondition}

$sp$ = most specific property which is guaranteed
after executing the code.

Basically tightest possible description of the
output state.

If we know a stronger fact than some $Q$, then

\[
sp \Rightarrow Q
\]

so strongest postcondition gives more information.

\vspace{8pt}

\textbf{WP}

Weakest precondition is kind of reverse direction.

Given the condition we want after the program,
find the least restrictive condition required before
the program.

\[
\{P\}\ C\ \{Q\}
\]

we calculate what $P$ has to be so that $Q$ is guaranteed.

Example:

\[
\{P\}\quad x:=x+5\quad\{x>15\}
\]

After assignment we need

\[
x+5>15
\]

therefore

\[
x>10
\]

hence

\[
WP(x:=x+5,\ x>15)=x>10
\]

WP is obtained by reasoning backwards from the desired
postcondition.

\[
P \Rightarrow WP
\]

\vspace{8pt}

Some basic WP rules:

\[
wp(skip,Q)=Q
\]

\[
wp(abort,Q)=false
\]

For assignment,

\[
wp(x:=E,Q)=Q[x\leftarrow E]
\]

i.e. substitute $E$ for $x$ in $Q$.

For sequential statements,

\[
wp(S_1;S_2,Q)
=
wp(S_1,wp(S_2,Q))
\]

For conditional:

\[
wp(if\ B\ then\ S_1\ else\ S_2,Q)
\]

\[
=
(B\Rightarrow wp(S_1,Q))
\land
(\neg B\Rightarrow wp(S_2,Q))
\]

\vspace{8pt}

\textbf{loop case}

Loops are harder because we can have many iterations.

Need some invariant $I$ which stays true while
the loop is running.

\vspace{8pt}

\textbf{invariant}

An invariant is a property which holds at a particular
control point in the program.

For a loop:

\[
\text{before loop}
\]

and at the beginning of every iteration,
the invariant should hold.

After the loop also it should still give us
something useful.

Three things to check:

\[
\text{initialisation}
\]

invariant is true before first iteration.

\[
P\Rightarrow I
\]

Then preservation:

\[
I\land B\land Body \Rightarrow I'
\]

one execution of the loop body should preserve $I$.

Then exit:

\[
I\land\neg B\Rightarrow Q
\]

when loop stops, invariant + exit condition should
give the required postcondition.

\vspace{8pt}

So roughly:

\[
P
\rightarrow I
\rightarrow I
\rightarrow I
\rightarrow \cdots
\]

and when the loop condition becomes false,

\[
I+\text{exit condition}\rightarrow Q
\]

\vspace{8pt}

\textbf{Dafny}

Dafny can use these conditions to automatically
check programs.

\begin{itemize}[leftmargin=18pt,itemsep=1pt]
    \item \texttt{requires} = precondition
    \item \texttt{ensures} = postcondition
\end{itemize}

For example,

\[
\texttt{requires } P
\]

means caller must establish $P$.

and

\[
\texttt{ensures } Q
\]

means Dafny tries to prove $Q$ after execution.

For some loops Dafny cannot finish the proof unless
we give it the right invariant.

\vspace{8pt}

\textbf{termination}

For loops that should terminate, need a quantity/relation
that is well-founded and keeps decreasing.

This is called a variant / decreases measure.

Idea:

\[
V_0 > V_1 > V_2 > \cdots
\]

and because the relation is well-founded, we cannot
keep decreasing forever.

\vspace{8pt}

\textbf{what Dafny does underneath}

Dafny converts the program into logical formulas.

These are called

\[
\text{Verification Conditions (VCs)}
\]

The VC is sent to an SMT solver, Z3.

The idea is to check the negation of the property
that Dafny wants to prove.

Then:

\[
\text{UNSAT}
\Rightarrow
\text{no counterexample}
\Rightarrow
\text{property is valid}
\]

while

\[
\text{SAT}
\Rightarrow
\text{counterexample exists}
\]

and

\[
\text{UNKNOWN}
\Rightarrow
\text{solver could not settle it}
\]

So overall:

\[
\boxed{\text{Dafny program}
\rightarrow
\text{VC}
\rightarrow
\text{Z3}
\rightarrow
\text{proof / counterexample}}
\]

\end{document}